% !TEX program = xelatex
% ============================================================================
% 中文课程作业 LaTeX 模板
% 中文课程作业 LaTeX 模板
% ============================================================================
% 使用方法：
%   1. 修改 \title、\author 中的个人信息
%   2. 在正文各 section 中填写内容
%   3. 如需图片，将图片文件放在 ../figures/ 目录下
%   4. 编译：latexmk -xelatex main.tex
% ============================================================================

\documentclass[12pt, a4paper]{ctexart}

% --- 页面与排版 ---
\usepackage{geometry}
\geometry{left=2.5cm, right=2.5cm, top=2.5cm, bottom=2.5cm, headheight=0pt, headsep=0pt}

\usepackage{setspace}
\onehalfspacing

% --- 数学与表格 ---
\usepackage{amsmath, amssymb}
\usepackage{booktabs}
\usepackage{graphicx}
\usepackage{hyperref}

% --- 图片路径：指向项目 figures/ 目录（可按需修改） ---
\graphicspath{{../outputs/figures/}}

% --- 标题格式 ---
\ctexset{
  section = {
    format    = \large\bfseries\raggedright,
    afterskip = 0.8em plus 0.2em minus 0.2em,
  },
  subsection = {
    format    = \normalsize\bfseries\raggedright,
    afterskip = 0.6em plus 0.1em minus 0.1em,
  },
}

% --- 页面风格：无页眉，仅页码 ---
\pagestyle{plain}

% --- 标题信息（修改这里） ---
\title{\vspace{-1.5em}作业标题\vspace{-0.5em}}
\author{姓名 \quad 邮箱 \quad 学院}
\date{}

\begin{document}

\maketitle
\vspace{-0.8em}

% --- 摘要 ---
\noindent\textbf{摘要：}%
在此处撰写摘要。摘要应简要概述研究问题、方法、主要结果和结论。篇幅控制在 150--300 字。

\vspace{1em}

% ============================================================
\section{引言}
% ============================================================

引言部分介绍研究背景、相关工作和研究动机。

% ============================================================
\section{方法}
% ============================================================

方法部分描述你的理论框架、模型设定和估计策略。对于包含数学公式的方法，可以使用行内公式 \(E = mc^2\) 或行间公式：

\begin{equation}
    \hat{\theta} = \frac{1}{n} \sum_{i=1}^{n} f(X_i)
    \label{eq:example}
\end{equation}

公式~\eqref{eq:example} 展示了某估计量的定义。

% ============================================================
\section{实验设置}
% ============================================================

\subsection{数据}

描述数据来源、样本量、变量定义和预处理步骤。

% --- 图片占位示例 ---
% 如需插入图片，请使用以下格式：
% \begin{figure}[t]
%     \centering
%     \includegraphics[width=0.80\linewidth]{your-figure.pdf}
%     \caption{图片标题}
%     \label{fig:example}
% \end{figure}
%
% 如果图片尚未准备好，可以使用占位框：
% \begin{figure}[t]
%     \centering
%     \fbox{\parbox{0.85\linewidth}{\centering 图片占位：\texttt{your-figure.pdf}}}
%     \caption{图片标题（占位）}
%     \label{fig:placeholder}
% \end{figure}

\subsection{实现细节}

描述编程语言、关键库、超参数设置等。

% ============================================================
\section{结果与讨论}
% ============================================================

结果与讨论部分呈现主要发现并分析其含义。

% ============================================================
\section{结论}
% ============================================================

结论部分总结主要发现，讨论局限性，指出未来方向。

% ============================================================
% 参考文献（如需使用 BibTeX，取消以下注释并创建 refs.bib）
% \bibliographystyle{plain}
% \bibliography{refs}
% ============================================================

\end{document}
